\documentclass{../../../unipd-notes}

\unipdsetup{
  course = {Fondamenti di Ingegneria Informatica},
  short-course = {Fondamenti di Ingegneria},
  document-type = {Appunti delle lezioni}
}

\begin{document}
\chapter{Comunicazione tra processi}
\section{Buffer produttore--consumatore}\label{sec:elements}

\begin{theorem}\label{thm:identity}
Se gli accessi al buffer sono protetti da mutua esclusione, il numero di elementi rimane compreso tra zero e la capacità $N$.
\end{theorem}

\begin{equation}\label{eq:identity}
  0 \leq n_{\mathrm{buffer}} \leq N
\end{equation}

\begin{figure}[htbp]
  \centering
  \begin{tikzpicture}[node distance=16mm]
    \node[draw=unipd-primary,rounded corners] (source) {Produttore};
    \node[draw=unipd-primary,rounded corners,right=of source] (target) {Consumatore};
    \draw[draw=unipd-accent,-{Latex},thick] (source) -- (target);
  \end{tikzpicture}
  \caption{Flusso dei dati attraverso il buffer condiviso}
  \label{fig:reference}
\end{figure}

\begin{table}[htbp]
  \centering
  \caption{Stato dei semafori del buffer}
  \label{tab:types}
  \begin{tabular}{ll}
    \toprule
    Semaforo & Valore iniziale \\
    \midrule
    \texttt{mutex} & $1$ \\
    \texttt{empty} & $N$ \\
    \texttt{full} & $0$ \\
    \bottomrule
  \end{tabular}
\end{table}

\begin{code}[language=Python,caption={Verifica della disponibilità del buffer},label={lst:identity}]
def can_consume(items):
    return items > 0
\end{code}

\begin{algorithm}[htbp]
  \caption{Prelievo di un elemento dal buffer}
  \label{alg:identity}
  \KwInput{Un buffer condiviso non vuoto}
  \KwOutput{L'elemento più vecchio}
  attendi(\texttt{full})\;
  attendi(\texttt{mutex})\;
  $x\gets$ rimuovi elemento\;
  segnala(\texttt{mutex})\;
  segnala(\texttt{empty})\;
  \Return{$x$}\;
\end{algorithm}

L'invariante del \cref{thm:identity}, espresso nell'\cref{eq:identity}, impedisce overflow e underflow. La \cref{fig:reference} riassume il flusso, mentre la \cref{tab:types}, il \cref{lst:identity} e l'\cref{alg:identity} descrivono inizializzazione e prelievo.
\end{document}
